\section{Introduction}
\label{sec:introduction}

This note continues the bachelor thesis~\cite{machacek2025method}. There the Galerkin matrix of the electric field
integral equation (\acs{EFIE}) with rooftop basis functions on a uniform rectangular grid is assembled from three Toeplitz
generators. Their entries are four-dimensional integrals of the Green's function
\begin{equation}
    \label{eq:green}
    g(R) = \frac{\mathrm{e}^{-\I k R}}{R}, \qquad R = \NORM{\V{r}-\V{r}'},
\end{equation}
over pairs of cells, weighted by the linear profiles of the rooftops (the factor $1/(4\uppi)$ is kept outside). They
are integrated with a tensor Gauss--Legendre (\GL{}) rule of order $q$, and the singularity is extracted on the
coincident cell only (Entry~2.1 in Section~2.6 and the appendix of~\cite{machacek2025method}).

Cells that share an edge or only a vertex meet $R=0$ as well, on the boundary of the integration domain. There the \GL{}
error does not decrease with the order: the generator entries keep a relative error of about $10^{-3}$ for cells of
$\lambda/10$ at every order between $3$ and $16$, and extraction on the coincident cell alone changes nothing measurable.
The note therefore extracts the non-smooth terms of~\eqref{eq:green},
\begin{equation*}
    \frac{1}{R},\qquad -\frac{k^2}{2}R,\qquad \frac{k^4}{24}R^3,
\end{equation*}
on all nine touching cells and gives the closed forms of every static integral this needs.

\paragraph{Three terms are enough.} With the three terms the entries reach $10^{-13}$--$10^{-14}$ at the order $12$--$16$,
and the computed current is at the rounding floor of double precision. A fourth term ($R^5$) would act below that floor,
so no further extraction is needed.

\paragraph{The way to the closed forms.} Section~\ref{sec:setting} expresses the three generators through seven
cell-pair moments, so that only integrals over pairs of cells with linear weights are left. Section~\ref{sec:extraction}
splits the Green's function into the three extracted terms and a smooth remainder; the remainder stays with the \GL{}
rule, and the extracted terms leave the static moments with the kernels $R^{-1}$, $R$ and $R^3$. Section~\ref{sec:reduction}
reduces these four-dimensional integrals to two dimensions, and Section~\ref{sec:antiderivatives} gives the
antiderivatives of the resulting integrands. Section~\ref{sec:touching} puts them together: every static moment of a
touching cell pair is a short sum of antiderivative values. The note is written as a manual to the implementation; all
closed forms are listed in Appendices~\ref{app:antiderivatives} and~\ref{app:explicit} in the form of the thesis
appendix.

\paragraph{Conventions.} The time convention is $\mathrm{e}^{\I\omega t}$, and $\arsinh x=\log\bigl(x+\sqrt{x^2+1}\bigr)$;
the logarithms of the thesis appendix are $\log\bigl((\sqrt{a^2+b^2}+a)/b\bigr)=\arsinh(a/b)$. The notation follows
the thesis~\cite{machacek2025method}.
